\subsection{Presentation of the Poincaré groups}

THEOREM 1. --- Let X be a compact and loop-disentanglable topological space, and let x be a point of X. The Poincaré group \(\pi_{1}(X,x)\) is of finite presentation.

The arc-connected component of x in X is open, closed, and loop-disentanglable; this allows us to assume that the space X is connected. Since the space X is loop-disentanglable, the X-space \(E = \lambda_{x}(X)\) , equipped with the endpoint map, is a nonempty and simply path-connected covering (IV, p. 342, th. 1). The group \(G = \text{Aut}_{\text{X}}(\text{E})\) is isomorphic to \(\pi_{1}(\text{X}, x)\) ; it is therefore a matter of proving that the group G is finitely presented.

Since X is compact, every point x of X has a compact neighborhood \(K_{x}\) over which the covering E is trivializable (TG, I, p. 65, corollary). Since X is locally connected, every \(x \in X\) has an open neighborhood \(W_{x}\) , connected and contained in \(K_{x}\) . Let F be a finite subset of X such that the \(W_{x}\) , for \(x \in F\) , cover X. Let n be the cardinality of F.

Let us show by induction that there exists, for every integer \(k\) such that \(1 \leqslant k \leqslant n\) , a subset A of cardinality \(k\) contained in F and, for \(x \in A\) , a section \(s_x\) of \(p\) over \(K_x\) , such that the union of the \(s_x(W_x)\) , for \(x \in A\) , is a connected subset of E. The assertion is true for \(k = 1\) . Suppose it is true for an integer \(k\) such that \(1 \leqslant k < n\) and let us show that it is true for \(k + 1\) . Let A be a subset of F of cardinality \(k\) and, for all \(x \in A\) , let \(s_x\) be a section of E over \(K_x\) such that \(\bigcup_{x \in A} s_x(W_x)\) is connected. The open sets \(\bigcup_{x \in A} W_x\) and \(\bigcup_{x \in F-A} W_x\) are nonempty and cover X; their intersection is therefore nonempty, since X is connected. There thus exist \(x \in A\) , \(y \in F - A\) and \(z \in W_x \cap W_y\) . Set \(A' = A \cup \{y\}\) and let us choose a section \(s_y\) of E over \(K_y\) such that \(s_y(z) = s_x(z)\) . The connected open sets \(\bigcup_{p \in A} s_p(W_p)\) and \(s_y(W_y)\) are nonempty and have a point in common; their union is therefore connected. This proves the assertion for \(k + 1\) . By induction, it is therefore true for every integer \(k \in \{1, \ldots, n\}\) .

Applying the above to k = n , there then exists, for every \(x \in F\) , a section \(s_x\) of E over \(K_x\) , such that \(U = \bigcup_{x \in F} s_x(W_x)\) is an open connected subset of E. We have \(p(U) = X\) . Since the group G acts transitively in each fiber of the covering E, we have GU = E.

The closure of U is contained in \(\bigcup_{x\in\mathrm{F}}s_x(\mathrm{K}_x)\) , therefore is compact. The action of G on E is proper (I, p. 96, cor. 1), hence the set of pairs \((g,x)\in\mathrm{G}\times\mathrm{E}\) such that \(x\in\overline{\mathrm{U}}\) and \(gx\in\overline{\mathrm{U}}\) is a compact subset of \(\mathrm{G}\times\mathrm{E}\) (TG, I, p. 77, prop. 6). It follows that the set of \(g\in\mathrm{G}\) such that \(\overline{\mathrm{U}}\cap g\overline{\mathrm{U}}\neq\varnothing\) is compact. Since it is discrete, it is finite. A fortiori, the set of \(g\in\mathrm{G}\) such that \(\mathrm{U}\cap g\mathrm{U}\neq\varnothing\) is finite. The theorem thus follows from prop. 10 of I, p. 136.

